% Online Resource 1 for the AC--MOT JRTIP manuscript.
% This file contains supporting evidence moved out of the 12-page main paper.
\documentclass[11pt]{article}
\usepackage[T1]{fontenc}
\usepackage{amsmath,amssymb}
\usepackage{float}
\newfloat{algorithm}{tbp}{loa}
\floatname{algorithm}{Algorithm}
\usepackage{algpseudocode}
\usepackage{booktabs}
\usepackage{longtable}
\usepackage{geometry}
\usepackage{enumitem}
\geometry{a4paper,margin=0.72in}
\setlength{\parindent}{0pt}
\setlength{\parskip}{4pt}
\renewcommand{\thetable}{S\arabic{table}}
\renewcommand{\thealgorithm}{S\arabic{algorithm}}
\newcommand{\checkpointsha}{\texttt{c6f98247\allowbreak{}e7c731a5\allowbreak{}67aaf37c\allowbreak{}a7b808be\allowbreak{}ff588b68\allowbreak{}7623e028\allowbreak{}b25a2e1d\allowbreak{}846e96e7}}

\title{Online Resource 1\\[2pt]
AC--MOT: Adaptive Detector Operating-Point Control for Real-Time Aerial Multi-Object Tracking}
\author{Ahmed Gouda Ismail, Mohamed S. Mohamed, and Tarek Ahmed Mahmoud}
\date{}

\begin{document}
\maketitle

\begin{center}
\textit{Journal of Real-Time Image Processing}\\
Computer Engineering and Artificial Intelligence Department, Military Technical College, Cairo, Egypt;\\
Computer Engineering Department, Egypt University of Informatics, Cairo, Egypt\\
Corresponding author: Ahmed Gouda Ismail, \texttt{a7medgouda1@gmail.com}
\end{center}

This supplementary document preserves supporting historical, implementation, and
provenance evidence that is not required to understand the primary contribution
or to verify the main-paper claims. Historical and modern protocols remain
separate. All values below are copied from the frozen manuscript evidence and
are not new experiments.

\section*{S1. Extended historical evidence}

\subsection*{S1.1 Component ablation}
The archived component progression is shown in full. The rows are historical
evidence and are not pooled with AC--MOT v3's results.

\begin{table}[ht]
\centering\small
\caption{Full historical component ablation.}
\label{tab:s1-ablation}
\begin{tabular}{lrrrrr}
\toprule
Configuration & MOTA & HOTA & IDF1 & IDS & FPS\\
\midrule
OLD--A0 & 17.633 & 29.837 & 30.892 & 283 & 44.085\\
OLD--A1 tuned tracker & 17.802 & 30.864 & 32.854 & 217 & 44.511\\
OLD--A2 confidence/NMS & 17.636 & 31.384 & 33.727 & 210 & 40.651\\
OLD--A2R resolution & 18.425 & 32.446 & 35.669 & 241 & 39.562\\
OLD--A3 full AC--MOT & 18.165 & 33.064 & 36.296 & 271 & 37.962\\
\bottomrule
\end{tabular}
\end{table}

\subsection*{S1.2 Historical VisDrone profiles}
The historical frozen comparison contains separate operating profiles. V1 is
quality-oriented, whereas V2 deliberately trades quality for identity and
runtime behaviour.

\begin{table}[ht]
\centering\small
\caption{Full historical VisDrone comparison.}
\label{tab:s1-visdrone}
\begin{tabular}{lrrrrr}
\toprule
System & MOTA & HOTA & IDF1 & IDS & FPS\\
\midrule
Baseline & 19.729 & 28.430 & 32.724 & 1235 & 36.528\\
Old AC--MOT & 23.236 & 32.698 & 39.516 & 1061 & 41.957\\
V1 Trial 24 & 26.948 & 33.835 & 41.546 & 1184 & 38.985\\
V2 Trial 22 & 23.792 & 31.218 & 37.870 & 919 & 46.024\\
\bottomrule
\end{tabular}
\end{table}

\subsection*{S1.3 Historical V1 uncertainty and UAVDT transfer}

\begin{table}[ht]
\centering\small
\caption{Full archived V1 paired bootstrap against the historical baseline.}
\label{tab:s1-bootstrap}
\begin{tabular}{lrr}
\toprule
Metric & Difference & 95\% CI\\
\midrule
MOTA & +7.2195 & [5.4362, 9.2934]\\
HOTA & +5.4051 & [4.1273, 6.9022]\\
IDF1 & +8.8220 & [6.9005, 11.0537]\\
IDS & -- & not significant\\
\bottomrule
\end{tabular}
\end{table}

\begin{table}[ht]
\centering\small
\caption{Full historical zero-retuning UAVDT transfer.}
\label{tab:s1-uavdt}
\begin{tabular}{lrrrr}
\toprule
System & MOTA & HOTA & IDF1 & IDS\\
\midrule
Baseline & 13.841 & 24.085 & 27.887 & 558\\
V1 & 17.399 & 28.390 & 34.453 & 321\\
V2 & 16.118 & 26.930 & 32.014 & 308\\
\bottomrule
\end{tabular}
\end{table}

\subsection*{S1.4 V1 optimization record}
The first formal optimization-derived configuration used fixed pretrained
YOLOv8n and tuned ByteTrack on VisDrone2019-MOT-val. Test data were not used
during search. The canonical Optuna/TPE stage searched five non-negative SCI
weights normalized to sum one, confidence and NMS endpoints selected from the
supported validation sets, and two sorted SCI thresholds. The preceding
validation screens supplied resolution levels 512, 928, and 960; the temporal
design supplied a seven-entry history and stride 10. The final directions were
maximize MOTA, minimize IDS, and maximize FPS, followed by the validation-only
feasibility rule FPS gate and IDS no greater than the OLD--A3 reference, then
highest MOTA with the documented tie breaks. Trial24 was frozen.

The exact V1 weights are crowd 0.1294928, tiny 0.2217477, edge 0.4337134,
night 0.0535577, and blur 0.1614885. The confidence endpoints are 0.30 and
0.40, both NMS endpoints are 0.35, and the SCI thresholds are 0.1353494 and
0.2872868. The V1 held-out bootstrap intervals and complete historical rows
are in Tables~\ref{tab:s1-visdrone} and~\ref{tab:s1-bootstrap}; these are
archived records rather than new experiments.

\subsection*{S1.5 V2 optimization record}
V2 reused the V1 validation design and changed the explicit objective to
maximize MOTA and minimize IDS with only FPS $\geq25$ as a feasibility gate.
The implementation constructs the feasible Pareto front by the standard
no-worse-in-both-objectives and strictly-better-in-one rule, then records
highest-MOTA, lowest-IDS, and balanced candidates. The balanced candidate uses
the documented equal-weight normalized MOTA/IDS score with HOTA, IDF1, FPS,
and trial-index tie breaks. Trial22 was the official selected configuration;
49 trials (0--48) completed in the archived study. Its validation record is
MOTA 19.330, HOTA 31.651, IDF1 34.226, IDS 168, and 51.406 FPS. The separate
post-selection testdev rerun is 23.792 MOTA, 31.218 HOTA, 37.870 IDF1, 919
IDS, and 46.024 FPS, and is not the selection record.

\subsection*{S1.6 Cross-pipeline audit}
The U2MOT host is the published ICCV implementation with YOLOX-X. V1 cue
weights and transforms were retained, while U2MOT validation distributions
provided unsupervised q33/q67 boundaries because the V1 boundaries collapsed
all frames into the hard region. The held-out custom evaluator reported
baseline/controller MOTA 53.9/53.7668, IDF1 69.8/69.8494, and IDS 1239/1152;
no HOTA, confidence interval, or matched speed claim is available.

The SparseTrack host was evaluated on MOT17 \texttt{val\_half}. Its transfer changes
only NMS using an edge cue sampled every ten frames, with a threshold learned
by leave-one-sequence-out development on the same seven sequences. Static and
adaptive MOTA are 76.8881 and 76.9252, with bootstrap CI
$[-0.041,+0.182]$ for the change; the manifest does not establish superiority.

\section*{S2. Extended modern result evidence}

\subsection*{S2.1 Full paired bootstrap metrics}
The main paper reports the headline HOTA, MOTA, IDF1, and IDS intervals.
The complete archived table additionally records false-positive and
false-negative changes.

\begin{table}[ht]
\centering\scriptsize
\caption{Complete modern paired bootstrap changes versus the declared baseline.}
\label{tab:s2-bootstrap}
\begin{tabular}{llrrr}
\toprule
Host & Metric & Difference & CI low & CI high\\
\midrule
ByteTrack & HOTA & +2.937 & +2.356 & +3.375\\
ByteTrack & MOTA & +17.995 & +13.541 & +23.282\\
ByteTrack & IDF1 & +5.360 & +4.490 & +6.178\\
ByteTrack & IDS & -325 & -498 & -164\\
ByteTrack & FP & -14382 & -19176 & -9620\\
ByteTrack & FN & +1781 & +288 & +3594.075\\
OATrack & HOTA & +1.325 & +0.200 & +2.264\\
OATrack & MOTA & +6.318 & +4.741 & +7.879\\
OATrack & IDF1 & +2.258 & +0.548 & +3.474\\
OATrack & IDS & -169 & -292 & -42\\
OATrack & FP & -3994 & -6190 & -2494.900\\
OATrack & FN & -375 & -1462.575 & +670\\
\bottomrule
\end{tabular}
\end{table}

The bootstrap uses 5000 paired sequence resamples with seed 42. It quantifies
uncertainty across the available sequences; it does not make frames or
detector stochasticity independent.

\subsection*{S2.2 State occupancy}

\begin{table}[ht]
\centering\small
\caption{AC--MOT v3 state occupancy on the second VisDrone validation read.}
\label{tab:s2-occupancy}
\begin{tabular}{lrrr}
\toprule
Host & LOW & MEDIUM & HIGH\\
\midrule
ByteTrack & 7.38\% & 30.15\% & 62.47\%\\
OATrack & 7.38\% & 33.31\% & 59.31\%\\
\bottomrule
\end{tabular}
\end{table}

Occupancy verifies that the controller uses multiple states and that the
trajectory differs by tracker host. It does not establish that exact switching
timing causes the quality difference.

\subsection*{S2.3 Full modern UAVDT transfer}

\begin{table}[ht]
\centering\scriptsize
\caption{Full modern UAVDT transfer under the frozen policy and no retuning.}
\label{tab:s2-uavdt}
\begin{tabular}{llrrrrrrrr}
\toprule
Host & Policy & HOTA & MOTA & IDF1 & IDS & FP & FN & Precision & Recall\\
\midrule
ByteTrack & baseline & 45.923 & -1.573 & 58.400 & 601 & 279616 & 66050 & 49.571 & 80.625\\
ByteTrack & AC--MOT v3 & 47.004 & 16.028 & 61.767 & 486 & 212077 & 73702 & 55.751 & 78.381\\
OATrack & baseline & 47.728 & 19.607 & 62.101 & 517 & 197937 & 75612 & 57.270 & 77.820\\
OATrack & AC--MOT v3 & 47.640 & 22.873 & 62.931 & 486 & 183669 & 78774 & 58.800 & 76.893\\
\bottomrule
\end{tabular}
\end{table}

The transfer contains 20 sequences, 16,592 frames, and 49,776 detector
forward passes. It is a single frozen-policy run per host without bootstrap or
cross-validation.

\section*{S3. Implementation traces}

\subsection*{S3.1 Preliminary heuristic prototype execution order}
The early heuristic prototype is retained here as provenance and reproducibility
evidence, not as the primary AC--MOT generation. The controller sits between a
frame stream and a fixed detector--tracker pair. It reads cheap image
statistics and boxes retained by the preceding tracker, computes a scene state,
maps that state to detector confidence, NMS, and resolution, runs the detector,
and stores the new tracker output for the next decision. It never reads future
frames, ground truth, or the current action's post-association output.

The historical cues are crowd count $C_t=\min(N_t/30,1)$, the fraction of
previous boxes below $32\times32$ pixels, normalized Canny edge density, a
brightness-below-80 night indicator, and a Laplacian-variance-below-180 blur
indicator. The historical SCI was
\begin{equation*}
S_t=0.30C_t+0.30T_t+0.20E_t+0.10N_t+0.05B_t,
\end{equation*}
whose coefficients intentionally sum to 0.95. The Smart Calibrator used
$c_t=0.245-0.050S_t$, applied a documented $-0.012$ correction for crowded,
tiny, or night scenes, and clipped confidence to $[0.19,0.28]$. It used
$\eta_t=0.490-0.050S_t$, applied a $-0.012$ blur correction, and clipped NMS
to $[0.40,0.52]$. The historical resolution actions were 640, 736, and 832
pixels, with analysis stride 10 and a seven-entry moving mean. These details
are preserved to prevent the heuristic prototype from being confused with V1.
The 832-pixel action was used when $S_t>0.60$ or the tiny-object ratio exceeded
0.50; the 736-pixel action was used when $S_t>0.35$ or the scene was labelled
crowded or tiny; otherwise the action was 640 pixels. At an analysis frame,
crowd and tiny-object values came from the previously stored tracker boxes,
whereas edge, brightness, and blur came from the current reduced image. New
boxes were stored only after detector execution and association, so the
controller did not create a within-frame feedback loop. An empty previous
output gave a zero tiny ratio and a crowd value that saturated at one after
thirty boxes. Scene labels were evaluated alongside the scalar index, and the
seven-entry mean reduced one-frame noise at the cost of response delay.

\begin{algorithm}[ht]
\caption{Preliminary heuristic prototype controller}
\label{alg:s3-original}
\begin{algorithmic}[1]
\State initialize seven-entry SCI history and previous tracker boxes
\For{each frame $t$}
\If{$t=1$ or $t\bmod 10=1$}
\State compute $C_t,T_t,E_t,N_t,B_t$ from current image and prior boxes
\State append $S_t$ to the history and average the available entries
\State apply Smart Calibrator confidence, NMS, and resolution rules
\EndIf
\State run detector at held operating point and update the fixed tracker
\State store tracker boxes for future analysis
\EndFor
\end{algorithmic}
\end{algorithm}

\subsection*{S3.2 Execution trace for AC--MOT v3}
At an analysis step, \texttt{analyze\_visual} computes reduced grayscale
statistics from the current image: Canny edge fraction, mean brightness, and
Laplacian variance. The SceneLayer reads \texttt{prev\_boxes}, reported by the
tracker on the last processed frame, and computes box count, crowd value,
box-area tiny fraction, and the decaying object-count peak. On non-analysis
frames the previous state and action remain unchanged.

The implementation then evaluates the labels in order. Brightness can label a
frame as night, blur can label it as blur, tiny-object fraction can label it as
tiny, and crowd or edge conditions can label it as crowded. The raw target uses
the modern thresholds and tiny guard. If the count falls below 0.40 of the
decayed peak after the peak has reached five objects, the bounded recovery probe
can request HIGH for up to three analysis steps.

If the raw target is lower than the persistent state, the hold rules are
checked: HIGH is held when SCI remains above 0.50 or tiny remains above 0.40,
and MEDIUM is held when SCI remains above 0.25. The 30-frame dwell gate then
prevents a state change too soon after the previous change. Only after these
guards does the adapter translate the state into the frozen
resolution/NMS/confidence tuple, run the detector, update the tracker, and store
its boxes for a later decision.

\subsection*{S3.3 Reproducibility checklist}
\begin{enumerate}[leftmargin=*]
\item Use the detector checkpoint \texttt{dronefreak} / \texttt{visdrone-yolo11m}.
\item Keep the eight-sequence calibration split separate from the seven-sequence
      modern evaluation read.
\item Preserve the crowd history length of seven and analysis stride of ten.
\item Apply the exact thresholds used by AC--MOT v3, including the tiny guard, recovery probe, hysteresis,
      30-frame dwell, and action tuples.
\item Keep detector-side confidence distinct from OATrack's tracker-side
      \texttt{min\_conf=0.40}.
\item Use the class-agnostic project evaluator and label it as such rather than
      as an official VisDrone leaderboard protocol.
\item Pair the baseline and AC--MOT v3 by sequence and use 5000 resamples with seed 42 for
      the modern bootstrap.
\item Do not retune on the seven-sequence read or on UAVDT.
\end{enumerate}

\section*{S4. Provenance and selection audit}

The modern freeze records the detector hash, controller thresholds, action map,
tracker descriptions, random seed, calibration fold membership, forward-pass
count, and source-file checksums. Stage 3A selected the cue subset
\texttt{['crowd']}; this describes the weighted SCI, not the removal of tiny,
edge, brightness, or blur from the surrounding SceneLayer.

The seven-sequence modern VisDrone evaluation is a documented second validation
read. It must not be described as an untouched held-out confirmation. The
archive also records that frozen Trial 7 is not the literal argmax of the
documented selection rule under either archived eligibility interpretation. The
manuscript reports this deviation and does not rerun the search.

The matched static control uses 1088 resolution, NMS IoU 0.45, and confidence
0.40. Its HOTA is approximately 44.530 for ByteTrack and 48.269 for OATrack,
compared with AC--MOT v3's values 44.523 and 48.095. The shuffled timing control has
development mean $\Delta$HOTA 2.446957 for the causal schedule and 2.617228 for
the same-action-distribution shuffled schedule. These controls shift the safe
interpretation toward operating-point calibration and observation-set shaping;
they do not establish that random timing is globally superior.

\section*{S5. Boundary of the supplementary evidence}

The supplementary material supports auditability and future reproduction. It
does not add a new claim of universal tracker independence, official leaderboard
equivalence, 30-FPS operation, or a timing-specific causal gain. The main paper
retains the negative attribution evidence, the second-read disclosure, the
Trial 7 selection deviation, the runtime harness caveat, and the mixed OATrack
UAVDT transfer because those facts define the supported claim boundary.

\end{document}
